(1) 空间位置是 \((R\sin\theta\cos\Omega t,R\sin\theta\sin\Omega t,-R\cos\theta)\)。因此 \(T=mR^2(\dot\theta^2+\Omega^2\sin^2\theta)/2\)，\(V=-mgR\cos\theta\)；Lagrange方程为 \(\ddot\theta=\sin\theta(\Omega^2\cos\theta-g/R)\)。

(2) 固定转速下 \(V_{\rm eff}=-mgR\cos\theta-m\Omega^2R^2\sin^2\theta/2\)。驻点有 \(\theta=0,\pi\)，以及当 \(\Omega^2R\ge g\) 时 \(\cos\theta_*=g/(\Omega^2R)\) 的两对称位置 \(\pm\theta_*\)。临界转速为 \(\sqrt{g/R}\)。

(3) 曲率除以 \(mR^2\) 是小振动频率平方。最低点有 \(\omega_0^2=g/R-\Omega^2\)；两非零平衡有 \(\omega_*^2=\Omega^2\sin^2\theta_*>0\)；顶端不稳。临界点二次曲率为零，必须保留有效势的四次项才能判定有限振幅非谐运动。

(4) Jacobi积分 \(J=mR^2\dot\theta^2/2+V_{\rm eff}\) 为常数。实验室能量 \(E=mR^2(\dot\theta^2+\Omega^2\sin^2\theta)/2-mgR\cos\theta\)。两者相差 \(mR^2\Omega^2\sin^2\theta\)，所以电机功率为 \(\dot E=mR^2\Omega^2\sin(2\theta)\dot\theta\)。
